\chapter{Traceability}
\ChapterLabel{Traceability}
\index{traceability}
\index{requirements traceability matrix}
\section{Chapter Responsibility}
\label{Section:TraceabilityResponsibility}
Maintain the canonical links from stakeholder need to release evidence.

\noindent\textbf{Expected result:} Visible coverage and identifiable gaps across the lifecycle.

\section{How to Use This Chapter}
Maintain the canonical mapping from stakeholder need or approved constraint through requirement, design decision, implementation, test, verification, validation, and release evidence. Review both missing backward links and missing forward links.
\section{Project Traceability Status}
\label{Table:TraceabilityMatrix}
There are no project stories, requirements, use cases, tests, or verification
and validation records yet. Consequently there is no populated product mapping
to assess. Traceability is structurally prepared, not demonstrated as product
coverage. The owned question register in
\hyperref[Table:InitializationQuestions]{the Development Plan} links each gap
to its authoritative question. Future artifacts require reciprocal links.

\section{Reserved Template Identifiers}
Template example IDs US-001--US-003; REQ-B-001; REQ-F-001--REQ-F-003;
REQ-Q-001; REQ-I-001; REQ-S-001; REQ-D-001; REQ-C-001; UC-001--UC-003;
TC-001--TC-009; VER-001--VER-009; VAL-001--VAL-003; and ARC-001 are reserved
and will not be reused for new project artifacts. They describe removed
instructional examples, not retired product requirements or implementation
evidence. Existing guide and diagram labels are preserved when their targets
remain in the derived book.
